\subsection{Quantum Gates}

\begin{definition}
A one-qubit quantum gate $U$ is a matrix $U \in \C^{2 \times 2}$ with:
\begin{enumerate}
    \item $U^\dagger U = UU^\dagger = I$;
    \item equivalently, $U = \ket{\phi_0}\bra0 + \ket{\phi_1}\bra1$, where $\{\ket{\phi_0}, \ket{\phi_1}\}$ forms an orthonormal basis.
\end{enumerate}
Common single-qubit gates are given in Table \ref{fig:single-qubit-gates}.
\end{definition}

\begin{figure}[H]
\centering
\renewcommand{\arraystretch}{1.1}
\begin{tabular}{c c c c}

\toprule
\textbf{Name} & \textbf{Circuit} & \textbf{Matrix} & \textbf{Bra-Ket} \\
\midrule

$X$ &
$\begin{quantikz}
\lstick{} & \targ{} & \qw
\end{quantikz}$ &
$\begin{pmatrix}
0 & 1 \\
1 & 0
\end{pmatrix}$ &
$\ket{1}\bra{0} + \ket{0}\bra{1}$ \\

$Y$ &
$\begin{quantikz}
\lstick{} & \gate{Y} & \qw
\end{quantikz}$ &
$\begin{pmatrix}
0 & -i \\
i & 0
\end{pmatrix}$ &
$i\ket{1}\bra{0} - i\ket{0}\bra{1}$ \\

$Z$ &
$\begin{quantikz}
\lstick{} & \gate{Z} & \qw
\end{quantikz}$ &
$\begin{pmatrix}
1 & 0 \\
0 & -1
\end{pmatrix}$ &
$\ket{0}\bra{0} - \ket{1}\bra{1}$ \\

$H$ &
$\begin{quantikz}
\lstick{} & \gate{H} & \qw
\end{quantikz}$ &
$\displaystyle\frac{1}{\sqrt{2}}
\begin{pmatrix}
1 & 1 \\
1 & -1
\end{pmatrix}$ &
$\ket{+}\bra{0} + \ket{-}\bra{1}$ \\

$S$ &
$\begin{quantikz}
\lstick{} & \gate{S} & \qw
\end{quantikz}$ &
$\begin{pmatrix}
1 & 0 \\
0 & i
\end{pmatrix}$ &
$\ket{0}\bra{0} + i\ket{1}\bra{1}$ \\

$T$ &
$\begin{quantikz}
\lstick{} & \gate{T} & \qw
\end{quantikz}$ &
$\begin{pmatrix}
1 & 0 \\
0 & e^{i\pi/4}
\end{pmatrix}$ &
$\ket{0}\bra{0} + e^{i\pi/4}\ket{1}\bra{1}$ \\

$R_x(\theta)$ &
$\begin{quantikz}
\lstick{} & \gate{R_x(\theta)} & \qw
\end{quantikz}$ &
$\begin{pmatrix}
\cos(\nicefrac{\theta}{2}) & -i\sin(\nicefrac{\theta}{2}) \\
-i\sin(\nicefrac{\theta}{2}) & \cos(\nicefrac{\theta}{2})
\end{pmatrix}$ &
$\begin{aligned}
\cos(\nicefrac{\theta}{2})\ket{0}\bra{0} 
&- i\sin(\nicefrac{\theta}{2})\ket{1}\bra{0} \\[-0.5ex]
- i\sin(\nicefrac{\theta}{2})\ket{0}\bra{1} 
&+ \cos(\nicefrac{\theta}{2})\ket{1}\bra{1}
\end{aligned}$ \\

$R_y(\theta)$ &
$\begin{quantikz}
\lstick{} & \gate{R_y(\theta)} & \qw
\end{quantikz}$ &
$\begin{pmatrix}
\cos(\nicefrac{\theta}{2}) & -\sin(\nicefrac{\theta}{2}) \\
\sin(\nicefrac{\theta}{2}) & \cos(\nicefrac{\theta}{2})
\end{pmatrix}$ &
$\begin{aligned}
\cos(\nicefrac{\theta}{2})\ket{0}\bra{0} 
&+ \sin(\nicefrac{\theta}{2})\ket{1}\bra{0} \\[-0.5ex]
- \sin(\nicefrac{\theta}{2})\ket{0}\bra{1}
&+ \cos(\nicefrac{\theta}{2})\ket{1}\bra{1}
\end{aligned}$ \\

$R_z(\theta)$ &
$\begin{quantikz}
\lstick{} & \gate{R_z(\theta)} & \qw
\end{quantikz}$ &
$\begin{pmatrix}
e^{-i\nicefrac{\theta}{2}} & 0 \\
0 & e^{i\nicefrac{\theta}{2}}
\end{pmatrix}$ &
$\begin{aligned}
e^{-i\nicefrac{\theta}{2}}\ket{0}\bra{0}
+ e^{i\nicefrac{\theta}{2}}\ket{1}\bra{1}
\end{aligned}$ \\

\bottomrule

\end{tabular}

\caption{Common single qubit quantum gates}
\label{fig:single-qubit-gates}
\end{figure}
\renewcommand{\arraystretch}{1}


\begin{theorem}[Euler decomposition]
\label{thm:euler-decomposition}
Every one-qubit unitary $U$ can be written as
\begin{align}
U = e^{i\alpha} R_Z(\beta) R_X(\gamma) R_Z(\delta),
\end{align}
where $\alpha,\beta,\delta \in [0,2\pi)$ and $\gamma \in [0,\pi]$.
\end{theorem}

\begin{proof}
This is a standard proof, seen for example in Section 4.2 of \cite{MikeAndIke}. Since $U$ is unitary, $|\det U|=1$. Choose $\alpha \in [0,2\pi)$ such that $\det U = e^{2i\alpha}$. Let $V := e^{-i\alpha}U$ satisfy $\det V=1$ while also being unitary. Now let
\begin{align}
V =
\begin{bmatrix}
a & b \\
c & d
\end{bmatrix}
\end{align}
Since the columns of $V$ are orthonormal, the second column must be a unit vector orthogonal to the first. Thus
\begin{align}
\begin{bmatrix}
b \\
d
\end{bmatrix}
&=
\lambda
\begin{bmatrix}
-\overline c \\
\overline a
\end{bmatrix}
\end{align}
for some $|\lambda|=1$. Hence
\begin{align}
V
&=
\begin{bmatrix}
a & -\lambda \overline c \\
c & \lambda \overline a
\end{bmatrix}.
\end{align}
Taking determinants gives
\begin{align}
1
=
\det V
&=
\lambda |a|^2 + \lambda |c|^2 \\
&=
\lambda (|a|^2+|c|^2) \\
&=
\lambda.
\end{align}
This also gives $|a|^2 + |c|^2 = 1$. Choose $\gamma,\eta,\mu \in \mathbb{R}$ such that
\begin{align}
a &= e^{i\eta}\cos(\gamma/2), \\
c &= e^{i\mu}\sin(\gamma/2).
\end{align}
Then
\begin{align}
V
&=
\begin{bmatrix}
e^{i\eta}\cos(\gamma/2) & -e^{-i\mu}\sin(\gamma/2) \\
e^{i\mu}\sin(\gamma/2) & e^{-i\eta}\cos(\gamma/2)
\end{bmatrix}.
\end{align}
On the other hand, matrix multiplication gives
\begin{align}
R_Z(\beta)R_X(\gamma)R_Z(\delta)
&=
\begin{bmatrix}
e^{-i(\beta+\delta)/2}\cos(\gamma/2)
&
-i e^{-i(\beta-\delta)/2}\sin(\gamma/2)
\\
-i e^{i(\beta-\delta)/2}\sin(\gamma/2)
&
e^{i(\beta+\delta)/2}\cos(\gamma/2)
\end{bmatrix}.
\end{align}
Thus it is enough to choose $\beta,\delta$ so that
\begin{align}
e^{i\eta} &= e^{-i(\beta+\delta)/2}, \\
e^{i\mu} &= -i e^{i(\beta-\delta)/2}.
\end{align}
Equivalently, modulo $2\pi$, solving gives
\begin{align}
\beta &= \mu-\eta+\frac{\pi}{2}, \\
\delta &= -\mu-\eta-\frac{\pi}{2}.
\end{align}
Therefore $V = R_Z(\beta)R_X(\gamma)R_Z(\delta)$. Since $U=e^{i\alpha}V$, we get
\begin{align}
U = e^{i\alpha}R_Z(\beta)R_X(\gamma)R_Z(\delta)
\end{align}
as desired.
\end{proof}

\begin{fact} 
\label{fact:tensor-properties}
For matrices $A,B,C,D$ over $\C$ of compatible dimensions, scalar $\lambda \in \C$, and vectors $\ket\phi$ and $\ket\psi$ over $\C$, the tensor product satisfies:
\begin{enumerate}
    \item Compatability with matrix multiplication: $(A \otimes B)(C \otimes D) = (AC) \otimes (BD)$, so $(A \otimes B)(\ket\phi \otimes \ket\psi) = (A\ket\phi) \otimes (B\ket\psi)$.
    \item Scalars commute: $\lambda(A \otimes B) = (\lambda A) \otimes B = A \otimes (\lambda B)$.
    \item Associativity: $(A \otimes B) \otimes C = A \otimes (B \otimes C)$.
    \item Unitary-preserving: If $A,B$ are unitary matrices, $A \otimes B$ is unitary.
\end{enumerate}
\end{fact}

\begin{definition}
    An \term{\(n\)-qubit gate} is a unitary matrix in
    \(\C^{2^n\times 2^n}\). If \(U_i\) is an \(n_i\)-qubit gate for
    \(i=1,\ldots,k\), and \(\sum_i n_i=n\), then
    \begin{align}
        U_1\otimes\cdots\otimes U_k
    \end{align}
    is an \(n\)-qubit gate by Fact~\ref{fact:tensor-properties}.
    Special named multi-qubit gates are given in Table~\ref{fig:multi-qubit-gates}.
\end{definition}

\begin{definition}
Many of the gates in Table \ref{fig:multi-qubit-gates} are \term{controlled} gates. Informally, controlled gates such as $\CNOT$ only apply their target operation when all controls
\begin{ZX}[circuit,thick lines]\rar & \zxCtrl{} \rar &\end{ZX}
are in the state $\ket{1}$. Formally, for any $n$-qubit gate $U$, the singly controlled gate, denoted $\cl U$, is defined by
\begin{align}
    \cl U := \ket0 \bra0 \otimes I + \ket1 \bra1 \otimes U.
\end{align}
This definition generalizes naturally to multiple controls. 
\end{definition}

By the above definition, $\CNOT$ is just $\cl X$, and the Toffoli gate is $\text{CC-}X$. Note that we may suppress the hyphen when the meaning is clear.

The gate $\cl Z$ is special because either qubit can be interpreted as the control: it only applies a phase when both qubits are in the state $\ket{1}$. To see this, note that
\begin{align}
    I = \ket0 \bra0 + \ket1 \bra1, \quad
    Z = I - 2\ket1\bra1.
\end{align}
Therefore,
\begin{align}
    \cl Z 
    &= \ket0 \bra0 \otimes I + \ket1 \bra1 \otimes (I - 2 \ket1 \bra1) \\
    &= I \otimes I - 2 \ket{11} \bra{11}.
\end{align}
Thus $\cl Z$ acts as the identity on every computational basis state except $\ket{11}$, where it applies a phase of $-1$.

\begin{figure}[H]
\centering
\renewcommand{\arraystretch}{1.1}
\begin{tabular}{c c c}

\toprule
\textbf{Name} & \textbf{Circuit} & \textbf{Bra-Ket} \\
\midrule

CNOT &
$\begin{quantikz}
\lstick{} & \ctrl{1} & \qw \\
\lstick{} & \targ{} & \qw
\end{quantikz}$ &
$\ket{0}\bra{0}\otimes I + \ket{1}\bra{1}\otimes X$ \\

CCNOT or Toffoli &
$\begin{quantikz}
\lstick{} & \ctrl{1} & \qw \\
\lstick{} & \ctrl{1} & \qw \\
\lstick{} & \targ{} & \qw
\end{quantikz}$ &
$\displaystyle \ket{11}\bra{11}\otimes X + \sum_{\substack{(a,b)\neq(1,1)\\a,b \in \{0,1\}}} \ket{ab}\bra{ab}\otimes I$ \\

SWAP &
$\begin{quantikz}
\lstick{} & \swap{1} & \qw \\
\lstick{} & \targX{} & \qw
\end{quantikz}$ &
$\displaystyle\sum_{i,j\in\{0,1\}} \ket{ij}\bra{ji}$ \\

C-$Z$ &
$\begin{quantikz}
\lstick{} & \ctrl{1} & \qw \\
\lstick{} & \ctrl{-1} & \qw
\end{quantikz}$ &
$\ket{0}\bra{0}\otimes I + \ket{1}\bra{1}\otimes Z$ \\

\bottomrule

\end{tabular}

\caption{Common multi-qubit quantum gates}
\label{fig:multi-qubit-gates}
\end{figure}
\renewcommand{\arraystretch}{1}

